\documentclass[11pt]{article}
\usepackage[utf8]{inputenc}
\usepackage{amsmath,amssymb}
\usepackage{booktabs}
\usepackage[margin=1in]{geometry}
\usepackage{setspace}

\newcommand{\grad}{\nabla}
\newcommand{\half}{\tfrac12}

\title{The Electron Lattice as a Dielectric:\\
       \large a periodic pinning potential for RealQM, and why it does not
       become a Frenkel--Kontorova model}
\author{Claes Johnson}
\date{PROSPECTUS --- not a draft. Revised \today.}

\begin{document}
\maketitle

\noindent\textbf{Status: a plan, with two measurements already made and one route already
closed.} The earlier version of this document proposed measuring the two Frenkel--Kontorova
coefficients and inferring the hopping barrier from them. One of the two came out cleanly; the
other does not exist in this construction, for a structural reason set out below. What remains is a
smaller and better-founded claim.

\section*{What is established}

\paragraph{A geometry in which the control is exact.} A rigid slide of a finite open chain is not a
symmetry operation --- it drags charge against the box and builds a dipole --- so $E(d)$ is not
periodic and no chain length repairs it. Lengthening the chain from seven cores to eleven made the
asymmetry \emph{worse}, not better. The repair is to make the translation cyclic: every electron
advances one site, the rightmost leaves, a new one enters at the left, and the configuration is
restored. No periodic Coulomb is needed for this, only a periodic \emph{labelling},
\[
\mathrm{dom}=\Bigl\lfloor\frac{x-x_0+a/2-d\,a}{a}\Bigr\rfloor+\mathrm{round}(d)
\quad(\mathrm{mod}\ N),
\]
where the added $\mathrm{round}(d)$ makes the label the index of the core the domain \emph{contains}
rather than a running counter --- without it the solver seeds each domain's density at the wrong
nucleus and the energies run below the physical floor.

\paragraph{The resulting scan, and its controls.} Seven cores at the calibrated $r_c=1.93$ and its
own equilibrium spacing $a=5.5$~a$_0$, the pattern slid through one full period:

\begin{center}
\small
\begin{tabular}{@{}rrrrrrrrr@{}}
\toprule
$d$ & $0.000$ & $0.125$ & $0.250$ & $0.375$ & $0.500$ & $0.625$ & $0.750$ & $1.000$\\
$E$ & $-1.63252$ & $-1.66402$ & $-1.71667$ & $-1.78418$ & $-1.84151$ & $-1.78416$ & $-1.71666$
    & $-1.63252$\\
\bottomrule
\end{tabular}
\end{center}

\noindent $E(1)=E(0)$ to every digit stored, and all three reflection pairs agree to
$1$--$2\times10^{-5}$~Ha. The noise floor is therefore about $0.5$~meV against an amplitude of
$5687$~meV: the control passes by four orders of magnitude, where the best attempt in the companion
article managed $24\times$ and six of eight failed outright.

\paragraph{The two numbers.} The minimum is at $d=\half$ --- walls on cores, the bond-centred
registry, independently reproducing the preference found by another route --- and the maxima are at
the commensurate registry:
\begin{align*}
\text{amplitude } V_0 &= E(0)-E(\half) = 0.20899~\mathrm{Ha} = 812~\mathrm{meV\ per\ electron},\\
K = \frac{E(0.375)-2E(0.5)+E(0.625)}{\delta u^2}
  &= 0.2426~\mathrm{Ha/a_0^2\ (total)} = 0.0347~\mathrm{Ha/a_0^2\ per\ electron}.
\end{align*}

\noindent $K>0$ is a restoring force: a field produces a bounded polarisation rather than a growing
displacement, so there is no DC current. \emph{This is an insulator established by a computed
response rather than by a failed barrier.} The polarisability follows, $\alpha=1/K\approx29$~a$_0^3$
per electron, between hydrogen's $4.5$ and lithium's $164$ --- the right neighbourhood for a valence
electron on a soft core, which is the only external check available.

\section*{What is closed, and why}

Frenkel--Kontorova needs a second coefficient, the stiffness $\kappa$ against non-uniform
displacement, after which $W=\sqrt{\kappa/K}$ would give the kink width and $E_{\rm
PN}\sim\exp(-cW/a)$ the barrier without ever differencing distant configurations. That was the point
of the exercise. It fails.

A sinusoidal displacement $u_m=A\sin(2\pi m/N)$, which closes periodically on the wrapped chain,
gives $\Delta E<0$ growing with $A$: $-0.0085$, $-0.0099$, $-0.0152$, $-0.0265$~Ha at
$A=0.25,\,0.30,\,0.35,\,0.40$. Fitting and subtracting the known substrate part leaves
$\kappa\approx-0.049$~Ha/a$_0^2$, so $W$ is imaginary and the model does not apply.

The cause is structural. In Frenkel--Kontorova $u_m$ is a particle \emph{position} and the spring
acts between particles. Here the particles are slabs, and position and width are not independent:
with walls at $w_m$, domain $m$ has centre $\half(w_{m-1}+w_m)$ and width $w_m-w_{m-1}$, so holding
the widths equal forces $w_m=ma+\mathrm{const}$ --- a rigid slide. \emph{Any non-uniform
displacement necessarily changes domain sizes}, and a wider domain lets its electron spread and
lower its kinetic energy. That gain --- the delocalisation the companion article measures at
$0.112$~Ha for a vacancy --- exceeds the Coulomb elastic cost.

So the negative $\kappa$ is not a chain that resists stretching in reverse. It is the system
preferring \textbf{unequal domains}, which also accounts for something already documented: free
boundaries never settle to a uniform partition, because the uniform partition is not stable.

\paragraph{Honest limits on that number.} The staircase noise is $\sim0.01$~Ha, so $\Delta E$ at the
largest amplitude is only $2.7$ times it, and the four ratios scatter between $-0.11$ and $-0.17$.
The \emph{sign} is consistent across all four amplitudes; the magnitude is one significant figure at
best. What is established is that $\kappa$ is not positive, hence that the mapping does not close ---
not the value of a negative $\kappa$.

\section*{What the paper would therefore claim}

Not a barrier, and not a Frenkel--Kontorova parameterisation. A \textbf{dielectric response computed
from Coulomb}: a periodic pinning potential of $812$~meV per electron with curvature
$K=0.0347$~Ha/a$_0^2$ and polarisability $29$~a$_0^3$, measured on a geometry whose translation
control is exact, together with the structural finding that the electron lattice is unstable to
non-uniform distortion and therefore admits no Frenkel--Kontorova description.

The competitive position is unchanged and should be stated plainly: DFT supplies dielectric
constants from first principles and is calibrated to do it. What is distinctive here is that the
pinning potential is a property of an explicit partition of space into unit charges, not of a smooth
density --- an interpretive difference, not an accuracy one.

\section*{Open, in order of value}

\begin{enumerate}\itemsep2pt
\item \textbf{Sub-cell interface resolution.} Everything above is limited by the staircase: the
partition assigns whole cells, so displacements below one cell flip no labels and $K$ cannot be
taken nearer the minimum than $u\approx0.34$~a$_0$. Two configurations $0.125$ and $0.250$ of a
spacing apart returned identical energies to five decimals --- the effect caught directly. This is
the one solver change that would improve every measurement in both papers.
\item \textbf{$K$ at a second spacing and a second $r_c$}, to see whether the dielectric response is
stable under the parameters, which no barrier ever was.
\item \textbf{The instability characterised.} At what wavelength is the uniform chain most unstable,
and does the distortion saturate? That is a Peierls-like question and the wrapped geometry can
now ask it properly.
\end{enumerate}

\end{document}
